\documentclass[a4paper,11pt,lualatex,ja=standard]{bxjsarticle}

% 数式
\usepackage{amsmath,amsfonts}
\usepackage{bm}
% 画像
\usepackage{graphicx}
% url
\usepackage{hyperref}

% ===== cover command =====
\newcommand{\cover}[3]{
  \thispagestyle{empty}
  \vspace*{4cm}
  \begin{center}
    {\LARGE #1}\\[2cm]
    {\Large #2}\\[3cm]
    {\large 氏名：#3}\\[1cm]
    {\large 学籍番号：s2510063}
  \end{center}
  \newpage
}
% =========================

\begin{document}

% ===== cover =====
\cover{I483}{report2}{小林智輝}
% =================

\title{レポートタイトル}
\author{氏名}
\date{\today}


\section{実行環境}

本課題では，MCUで取得したセンサーデータをMQTTで公開し，PC側でMQTTおよびKafkaを用いてデータの受信，処理，公開を行った。使用した実行環境を以下に示す。

\begin{itemize}
\item MCU：Raspberry Pi Pico W
\item MCU側開発環境：Thonny
\item MCU側言語：MicroPython
\item PC側開発環境：Visual Studio Code
\item PC側言語：Python 3.12.10
\item 通信方式：MQTT，Kafka
\item MQTTブローカ：150.65.230.59
\item MQTTポート：1883
\item 学生番号：s2510063
\item 使用ライブラリ：umqtt.simple，paho-mqtt，kafka-python
\end{itemize}

また，本課題で使用したセンサーを以下に示す。

\begin{itemize}
\item SCD41：CO2，温度，湿度
\item BH1750：照度
\item RPR-0521RS：照度，赤外線照度
\item DPS310：気圧，温度
\end{itemize}

\section{開発内容}

\subsection{MQTTによるセンサーデータの公開}

課題1(a)では，MCUで取得したセンサーデータをMQTTで公開した。MQTTトピックは，課題で指定された以下の形式に従った。

\begin{verbatim}
i483/sensors/STUDENT/SENSOR/DATA_TYPE
\end{verbatim}

ここで，\texttt{STUDENT}には自身の学生番号である\texttt{s2510063}を用いた。また，データは文字列形式で送信し，小数第2位までに整形した。

実行時の出力例を以下に示す。

\begin{verbatim}
publish: b'i483/sensors/s2510063/SCD41/co2' 1316.00
publish: b'i483/sensors/s2510063/SCD41/temperature' 26.34
publish: b'i483/sensors/s2510063/SCD41/humidity' 42.65
publish: b'i483/sensors/s2510063/BH1750/illumination' 586.67
publish: b'i483/sensors/s2510063/RPR0521/illumination' 853.65
publish: b'i483/sensors/s2510063/RPR0521/infrared_illumination' 99.00
publish: b'i483/sensors/s2510063/DPS310/air_pressure' 998.30
publish: b'i483/sensors/s2510063/DPS310/temperature' 28.55
\end{verbatim}

この結果から，SCD41，BH1750，RPR-0521RS，DPS310で取得した各データを，指定されたMQTTトピックへ公開できていることを確認した。

\subsection{MQTTによるデータ受信}

課題1(b)では，課題1(a)で公開されたデータをMQTTで受信するPC側プログラムを作成した。購読トピックには以下を指定した。

\begin{verbatim}
i483/sensors/s2510063/+/+
\end{verbatim}

この指定により，自身の学生番号に対応する全センサーのデータをまとめて受信できるようにした。

実行時の出力例を以下に示す。

\begin{verbatim}
MQTT connected: 0
Subscribe topic: i483/sensors/s2510063/+/+
2026-05-31 10:49:41 i483/sensors/s2510063/SCD41/co2 1146.00
2026-05-31 10:49:41 i483/sensors/s2510063/SCD41/temperature 27.68
2026-05-31 10:49:41 i483/sensors/s2510063/SCD41/humidity 39.69
2026-05-31 10:49:41 i483/sensors/s2510063/BH1750/illumination 271.67
2026-05-31 10:49:41 i483/sensors/s2510063/RPR0521/illumination 6.53
2026-05-31 10:49:41 i483/sensors/s2510063/RPR0521/infrared_illumination 1.00
2026-05-31 10:49:41 i483/sensors/s2510063/DPS310/air_pressure 998.05
2026-05-31 10:49:41 i483/sensors/s2510063/DPS310/temperature 29.90
\end{verbatim}

出力結果より，MQTTブローカへの接続に成功し，複数のセンサーデータを継続的に受信できていることを確認した。

\subsection{Kafkaによるデータ処理と公開}

課題1(c)では，Kafkaを利用して，BH1750の照度データとSCD41のCO2データを受信し，処理結果を別のKafkaトピックへ公開するプログラムを作成した。

購読したKafkaトピックを以下に示す。

\begin{verbatim}
i483-sensors-s2510063-BH1750-illumination
i483-sensors-s2510063-BH1750-temperature
i483-sensors-s2510063-SCD41-co2
\end{verbatim}

本プログラムでは，以下の2つの処理を行った。

\begin{itemize}
\item BH1750の照度データについて，最近5分間の値を用いてRolling Averageを計算する。
\item SCD41のCO2濃度が700ppmを超えた場合，Threshold Detectionとして\texttt{yes}を公開する。
\end{itemize}

実行時の出力例を以下に示す。

\begin{verbatim}
Kafka connected
Subscribe:
i483-sensors-s2510063-BH1750-illumination
i483-sensors-s2510063-BH1750-temperature
i483-sensors-s2510063-SCD41-co2

receive: i483-sensors-s2510063-SCD41-co2 1166.00
publish: i483-actuators-s2510063-co2_threshold-crossed yes

receive: i483-sensors-s2510063-BH1750-illumination 432.50
receive: i483-sensors-s2510063-BH1750-illumination 416.67
publish: i483-sensors-s2510063-BH1750_avg-illumination 424.59
\end{verbatim}

この結果から，Kafkaを用いてセンサーデータを受信し，Rolling AverageおよびThreshold Detectionの処理結果を，指定されたKafkaトピックへ公開できていることを確認した。

\subsection{Threshold Detectionに基づくLED点滅}

課題1(d)では，課題1(c)で行ったThreshold Detectionの結果に基づいて，MCUのLEDを点滅させた。

CO2濃度が700ppmを超えた場合，Kafkaトピック

\begin{verbatim}
i483-actuators-s2510063-co2_threshold-crossed
\end{verbatim}

に\texttt{yes}を公開する。その後，この情報はMQTTトピック

\begin{verbatim}
i483/actuators/s2510063/co2_threshold_crossed
\end{verbatim}

へ変換される。MCU側ではこのMQTTトピックを購読し，受信したメッセージが\texttt{yes}であった場合にLEDを点滅させるようにした。

実行時の出力例を以下に示す。

\begin{verbatim}
Kafka connected
Subscribe:
i483-sensors-s2510063-BH1750-illumination
i483-sensors-s2510063-BH1750-temperature
i483-sensors-s2510063-SCD41-co2
-------------------------------

receive: i483-sensors-s2510063-SCD41-co2 1132.00
publish: i483-actuators-s2510063-co2_threshold-crossed yes
LED should blink

receive: i483-sensors-s2510063-BH1750-illumination 400.00
receive: i483-sensors-s2510063-SCD41-co2 1116.00
receive: i483-sensors-s2510063-BH1750-illumination 399.17
publish: i483-sensors-s2510063-BH1750_avg-illumination 399.59
\end{verbatim}

この結果から，CO2濃度が700ppmを超えたことを検出し，Threshold Detectionの結果として\texttt{yes}を公開できていることを確認した。また，この結果をMCU側で受信し，LEDを点滅させる処理を実装した。

\section{運用}

課題2(a)では，課題1(a)で作成したMQTTによるセンサーデータ公開プログラムと，課題1(c)で作成したKafkaによる処理プログラムを継続的に運用した。

MCU側では，センサーから15秒ごとにデータを取得し，MQTTブローカへ公開した。PC側では，Kafkaを用いてデータを受信し，Rolling AverageおよびThreshold Detectionの処理を継続して行った。

運用中，以下の処理が継続的に行われていることを確認した。

\begin{itemize}
\item MQTTによるセンサーデータの公開
\item MQTTによるセンサーデータの受信
\item Kafkaによる照度データおよびCO2データの受信
\item BH1750照度データのRolling Averageの計算と公開
\item SCD41のCO2濃度に対するThreshold Detectionの実行
\item Threshold Detectionの結果に基づくMCUのLED点滅
\end{itemize}

以上より，MQTTとKafkaを用いたPub/Sub型通信により，センサーデータの収集，処理，およびアクチュエータ制御を継続的に実行できることを確認した。

\section{まとめ}

本課題では，MCUで取得したセンサーデータをMQTTで公開し，PC側でMQTTおよびKafkaを用いてデータの受信，処理，公開を行った。MQTTでは，指定されたトピック形式に従って複数のセンサーデータを公開し，PC側で受信できることを確認した。

また，Kafkaでは，BH1750の照度データに対するRolling Averageと，SCD41のCO2濃度に対するThreshold Detectionを実装した。さらに，Threshold Detectionの結果をMCU側で受信し，LEDを点滅させることで，Pub/Sub型通信を用いたセンサーデータ処理とアクチュエータ制御の一連の流れを確認できた。

\end{document}
